\paragraph{Graphs}
A graph $G=(V,E)$ is a structure composed of a set of vertices $V$ and a set edges $E$ which are unordered pairs of vertices. A degree of $v\in V$ $\deg(v)=|\set{u\in V|(u,v)\in E}|$ is the amount of edges from $v$.
A full graph, where $E=\set{(v,u)|v,u\in V}$ contains all possible edges is also called a clique. An anticlique is a graph where $E=\emptyset$ is empty. A bipartite graph $G=(U,V,E)$ is comprised of two anticliques, ie $E\cap V\times V=\emptyset$, $E\cap U\times U=\emptyset$.
A matching $M\subseteq E$ is a subset of edges that have no common endpoints $\not\exists e_1,e_2:\exists v\in V: v\in e_1\land v\in e_2$. A matching $M$ is maximum when it isn't a subset of another matching. A matching is maximal if there is no larger matching. A perfect matching $M^\star$ is a matching with $\frac{|V|}{2}$ edges, ie covering all vertices. We write $M(v)=u$ to denote that $(u,v)\in M$. We write $u$ is matched to $v$ to denote the same.
\paragraph{Online Algorithms}
An online algorithm is an algorithm in which the data isn't entirely known, but rather arrives over time and the algorithm makes decisions (usually irrevocably) without the knowledge of the remaining data. An online algorithm can be considered a greedy optimization algorithm. Compttitiveness is a way of measuring the performance of online algorithms and is a lower/upper bound (depending on whether the goal is to maximize/minimize a parameter). In the case of online algorithms for bipartite matching competitiveness of algorithm $\mathrm{ALG}$ is $\underset{G\text{ graph}}{\inf}\ \underset{\text{arrival order}}{\inf}\frac{\mathrm{ALG}(G)}{\mathrm{OPT}(G)}$ a guarantee on the minimum size of the matching found by $\mathrm{ALG}$ as a fraction of the largest matching $\mathrm{OPT}$. A $\frac{\mathrm{ONLINE}}{\mathrm{OFFLINE}}$ notation is sometimes used as an offline algorithm can find the maximal matching. In the case that an algorithm is randomized, it's competitiveness is $\underset{G\text{ graph}}{\inf}\ \underset{\text{arrival order}}{\inf}\frac{\E[\mathrm{ALG}(G)]}{\mathrm{OPT}(G)}$. An algorithm's performance on a graph is the infimum over data order of the expected size of found matching. An algorithm has an asymptotic performance is $\lim_{n\to\infty}$ of performance on graphs of size $n$
\paragraph{Linear Programming}
Linear Programming is the optimization of a linear function $f(x)=\sum\limits c_ix_i=c^Tx$ under linear constraints $Ax\leq b, x\geq 0$. For vectors of equal length $x,y$ we write $x\leq y$ when $\forall_i x_i\leq y_i$. Every linear program (the primal program) has a dual program associated with it, defined as $b^Ty$ under $A^Ty\geq c, y\geq 0$. If the primal program is a minimization of $f$ then the dual is a maximization and vice versa. A feasible solution is a value of $x$ that satisfies program constraints. An important propery of the dual program is Weak Duality: for all feasible primal, dual solutions $x,y$, $c^Tx\leq b^T y$. This allows for the dual program to be used as a bound for the primal. This technique can be used to bound competitiveness of algorithms if the algorithms can be expressed as a linear program